\section{Background}
\label{sec:background}

\paragraph{Prefix reuse.}
A server that has processed a prompt keeps the resulting model state. When the
next prompt starts with the same tokens, it evaluates only the new suffix. In
llama.cpp this happens implicitly within a slot (\texttt{cache\_prompt}), and
explicitly through endpoints that save a slot's state to a file and restore
it~\cite{llamacpp-server}. We use explicit save and restore, because it lets us
force each condition and then prove from the server's own counters which one
ran. We call the number of reused tokens the \emph{reuse boundary} $k$.

\paragraph{Micro-batches.}
llama.cpp evaluates a prompt in physical micro-batches of \nub{} tokens
(512 in this study). A prompt of length $n$ is split into
$[0,U), [U,2U), \dots$ with $U = \nub$.

\paragraph{What the state is.}
The three models differ in what a cached prefix consists of.
Qwen3-4B~\cite{qwen3} is a dense transformer: the state is the keys and values
of every prefix position in every layer.
Spark-X2.5-4B~\cite{spark} interleaves three sliding-window attention layers
with one full-attention layer; according to the model file's metadata the
window is 512 tokens, so the windowed layers need only the most recent 512
positions.
Ornith-1.5-9B~\cite{ornith} is a hybrid: its model file declares the
\texttt{qwen35} architecture, with full attention in every fourth layer and
linear-attention layers in between that carry a fixed-size recurrent state.
The model card does not describe this structure and calls the model dense; we
rely on the file metadata (\texttt{full\_attention\_interval} of 4 and a
recurrent-state block). A
recurrent state summarises the whole prefix, so an error in it is carried
forward to every later token, where a per-position cache keeps errors local to
the positions that produced them.
